% ============================================================
%  Black hole shadows in Topology-Generated Potential
%  Mateusz Serafin — April 2026
%  Self-contained paper for arXiv submission
% ============================================================
\documentclass[11pt,a4paper]{article}

\usepackage[utf8]{inputenc}
\usepackage[T1]{fontenc}
\usepackage{amsmath,amssymb,amsfonts,amsthm}
\usepackage{mathtools}
\usepackage[hidelinks]{hyperref}
\usepackage{geometry}
\usepackage{booktabs,array}
\usepackage[shortlabels]{enumitem}
\usepackage{xcolor}
\usepackage{graphicx}
\geometry{margin=2.3cm}

% --- Theorem environments ---
\newtheorem{theorem}{Theorem}[section]
\newtheorem{proposition}[theorem]{Proposition}
\newtheorem{lemma}[theorem]{Lemma}
\newtheorem{corollary}[theorem]{Corollary}
\theoremstyle{definition}
\newtheorem{definition}[theorem]{Definition}
\theoremstyle{remark}
\newtheorem{remark}[theorem]{Remark}

% --- Shortcuts ---
\newcommand{\PhiZero}{\varPhi_0}
\newcommand{\lP}{\ell_{\mathrm{P}}}
\newcommand{\rs}{r_{\mathrm{s}}}
\newcommand{\rph}{r_{\mathrm{ph}}}
\newcommand{\bcrit}{b_{\mathrm{crit}}}
\newcommand{\dd}{\mathrm{d}}
\renewcommand{\Phi}{\varPhi}

\begin{document}

% ============================================================
\title{Black hole shadows in Topology-Generated Potential:\\
  photon sphere shift and falsifiable predictions\\
  for the Event Horizon Telescope}
\author{Mateusz Serafin\thanks{%
  E-mail: \texttt{mateusz.serafin@tgp-project.org}.
  Repository: \url{https://github.com/MateuszSerafin/TGP}.}}
\date{April 2026}
\maketitle

% ============================================================
\begin{abstract}
Topology-Generated Potential (TGP) is a scalar-field framework in which
the physical constants $c$, $\hbar$, and $G$ are field-dependent while
the Planck length $\lP$ remains invariant.  The weak-field, static,
spherically symmetric metric takes the exponential (Yilmaz--Rosen) form
$\psi = \exp(-\rs/r)$, which reproduces all post-Newtonian predictions
of general relativity (GR) with $\gamma = \beta = 1$.

We present three levels of analysis.  \textbf{(1)~Exponential
extrapolation:} if taken to hold at all radii, the exponential metric
places the photon sphere at $R_{\mathrm{ph}} = \rs/2$ in curvature
coordinates, yielding a shadow $\sim 48\%$ smaller than Schwarzschild.
\textbf{(2)~Consistency check:} we prove analytically that the
exponential metric does \emph{not} satisfy the full nonlinear TGP field
equation --- the residual is $\mathcal{O}(\beta)$ at all radii ---
and that in isotropic coordinates the function $r\cdot\psi(r)$ is
monotonically increasing, so the exponential metric possesses
\emph{no photon sphere} in coordinate-invariant terms.  The
``$R_{\mathrm{ph}} = \rs/2$'' result is an artifact of applying the
curvature-coordinate photon-sphere formula to a metric form that is
not the true strong-field solution.  \textbf{(3)~Strong-field program:}
the TGP ontology requires $\Phi \to \infty$ at the center (frozen
interior), implying $\psi \to \infty$ and $g_{tt} \to 0$ --- an
effective horizon.  We show that \emph{any} profile with $\psi(0) =
\infty$ and $\psi(\infty) = 1$ must possess a photon sphere by the
intermediate value theorem, and that the shadow size depends on the
full nonlinear profile $\psi(r)$, which remains an open problem.

Black hole thermodynamics is preserved exactly
($S_{\mathrm{BH}} = k_B A / 4\lP^2$) and the no-singularity theorem
($K \to 0$) is robust.  Ringdown frequency shifts
($\delta f_0 \approx +0.6\%$, $\delta f_2 \approx -0.2\%$) provide
independent gravitational-wave tests.
\end{abstract}

\tableofcontents

% ============================================================
\section{Introduction}\label{sec:intro}

The direct imaging of the supermassive black hole in M87 by the Event
Horizon Telescope (EHT) collaboration~\cite{EHT-M87-2019,EHT-M87-VI}
and the subsequent image of Sagittarius~A*~\cite{EHT-SgrA-2022} have
opened an unprecedented window on strong-field gravity.  The bright
emission ring observed around each compact object is interpreted as the
lensed image of accreting plasma, with its diameter set by the \emph{%
critical impact parameter} $\bcrit$ --- the closest approach at which a
photon can still escape to infinity.  In general relativity, for a
Schwarzschild black hole of mass~$M$,
\begin{equation}\label{eq:bcrit-GR}
  \bcrit^{\mathrm{GR}}
  = 3\sqrt{3}\,\frac{GM}{c_0^2}
  \approx 2.598\,\rs,
  \qquad
  \rs \equiv \frac{2GM}{c_0^2}.
\end{equation}
The EHT measurements of M87* yield a shadow angular diameter of
$42 \pm 3\;\mu\mathrm{as}$~\cite{EHT-M87-VI}, consistent with the GR
prediction to within roughly $10\%$.  For Sgr~A*, the measured diameter
is $48.7 \pm 7\;\mu\mathrm{as}$~\cite{EHT-SgrA-2022}, again broadly
consistent with Schwarzschild.

These observations already constrain a number of modified gravity
theories and parametrized deviations from the Kerr
metric~\cite{Psaltis2020,Vagnozzi2023,Johannsen2013}.  At the same
time, the uncertainties --- dominated by source morphology and
interstellar scattering rather than by instrumental noise --- leave room
for $\mathcal{O}(10\text{--}30\%)$ deviations in the shadow
size~\cite{Perlick2022,Cunha2018}.  Any theory that makes a sharp,
quantitative prediction for the shadow diameter is therefore immediately
testable.

\medskip

Topology-Generated Potential (TGP) is a framework in which spacetime
itself is an emergent quantity generated by a scalar substrate
field~$\Phi$~\cite{Serafin2026-full,TGP-repo}.  The fundamental
constants $c$, $\hbar$, and~$G$ become field-dependent:
\begin{equation}\label{eq:scalings}
  c(\Phi) = c_0\!\left(\frac{\PhiZero}{\Phi}\right)^{\!1/2}\!,
  \quad
  \hbar(\Phi) = \hbar_0\!\left(\frac{\PhiZero}{\Phi}\right)^{\!1/2}\!,
  \quad
  G(\Phi) = G_0\,\frac{\PhiZero}{\Phi},
\end{equation}
where $c_0$, $\hbar_0$, $G_0$ are the vacuum values measured at spatial
infinity ($\Phi = \PhiZero$).  A key structural result is that the Planck
length $\lP = \sqrt{\hbar G/c^3}$ is \emph{invariant} under these
scalings, ensuring that quantum-gravitational physics remains universal.

In the weak-field regime, TGP reproduces the Schwarzschild solution to
all post-Newtonian (PPN) orders with $\gamma = \beta = 1$, matching every
Solar-System and binary-pulsar test catalogued by
Will~\cite{Will2014}.  The speed-of-gravity constraint from
GW170817/GRB~170817A~\cite{Baker2017} is likewise satisfied.  In the
strong-field regime, however, the exponential form of the TGP metric
produces qualitative departures: there is no true event horizon, no
curvature singularity, and the photon sphere shifts inward.

The purpose of this paper is to work out the observational consequences
of the TGP metric for the black hole shadow.  We derive the photon
sphere radius and critical impact parameter, tabulate predictions for
M87* and Sgr~A*, discuss the no-singularity theorem and thermodynamic
consistency, and assess the current observational status with explicit
caveats.

\medskip\noindent\textbf{Plan.}
Section~\ref{sec:metric} reviews the TGP exponential metric.
Section~\ref{sec:photon} derives the photon orbits.
Section~\ref{sec:no-sing} presents the no-singularity theorem.
Section~\ref{sec:thermo} verifies black hole thermodynamics.
Section~\ref{sec:obs} tabulates observational predictions.
Sections~\ref{sec:discussion} and~\ref{sec:conclusion} discuss
limitations and conclude.

% ============================================================
\section{TGP exponential metric}\label{sec:metric}

\subsection{Field equation and static solution}

The TGP substrate field $\Phi(\mathbf{x},t)$ satisfies a nonlinear wave
equation on flat background whose detailed derivation is given
in~\cite{Serafin2026-full,Serafin2026-nbody}.  For the purposes of this
paper, we need only the static, spherically symmetric vacuum solution.

In the weak-field (Newtonian) limit, $\Phi \approx \PhiZero$ and the
gravitational potential is $U = -GM/r$.  The TGP line element for a
static, spherically symmetric source takes the \emph{exponential metric}
form:
\begin{equation}\label{eq:exp-metric}
  \boxed{
    \dd s^2
    = -\frac{c_0^2}{\psi}\,\dd t^2
    + \psi\,(\dd r^2 + r^2\,\dd\Omega^2),
  }
\end{equation}
where
\begin{equation}\label{eq:psi-def}
  \psi(r) \equiv \exp\!\left(\frac{2U(r)}{c_0^2}\right)
  = \exp\!\left(-\frac{\rs}{r}\right),
  \qquad
  \rs = \frac{2GM}{c_0^2}.
\end{equation}
Here $\dd\Omega^2 = \dd\theta^2 + \sin^2\!\theta\,\dd\varphi^2$ is the
standard unit-sphere element.  The isotropic radial coordinate~$r$ is
used throughout.

\subsection{Key properties}

\begin{enumerate}[(i)]
\item \textbf{No coordinate singularity.}\quad
  Since $\psi = e^{-\rs/r} > 0$ for all $r > 0$, the metric
  coefficients are everywhere finite and positive.  There is no radius at
  which $g_{tt} = 0$ or $g_{rr} = \infty$; in particular, the surface
  $r = \rs$ is \emph{not} a horizon in the Schwarzschild sense.  The
  light cones pinch continuously but never close.

\item \textbf{Weak-field agreement with GR.}\quad
  Expanding $\psi$ to any finite order in $\rs/r$ and comparing with
  the PPN-expanded Schwarzschild metric in isotropic coordinates, one
  finds exact agreement at every order:
  \begin{equation}
    g_{tt} = -c_0^2 e^{+\rs/r}
    = -c_0^2\!\left(1 + \frac{\rs}{r} + \frac{1}{2}\frac{\rs^2}{r^2}
      + \cdots\right),
  \end{equation}
  \begin{equation}
    g_{rr} = e^{-\rs/r}
    = 1 - \frac{\rs}{r} + \frac{1}{2}\frac{\rs^2}{r^2} - \cdots\;.
  \end{equation}
  In particular, the PPN parameters are $\gamma = \beta = 1$, identical
  to GR~\cite{Will2014}.

\item \textbf{Strong-field departure.}\quad
  The exponential and Schwarzschild metrics agree perturbatively but
  differ non-perturbatively.  The Schwarzschild metric (in isotropic
  coordinates) has $g_{tt} = 0$ at $r = \rs/4$ (the horizon in isotropic
  form), while the TGP metric has $g_{tt} \to -c_0^2$ as $r \to 0$
  after passing through a minimum.  This is the root of the shadow-size
  discrepancy derived below.
\end{enumerate}

\subsection{Relation to the substrate field}

The conformal factor $\psi$ is related to the substrate field by
$\psi = \Phi/\PhiZero = \chi$.  In the weak-field exterior,
$\chi \approx 1 - \rs/(2r) + \cdots$, while in the strong-field
interior $\chi \gg 1$ (approaching the ``frozen'' phase described in
Section~\ref{sec:no-sing}).  The dimensionless field $\chi$ is the
natural variable for discussing both the duality structure and the
curvature invariants.

% ============================================================
\section{Photon orbits in the TGP metric}\label{sec:photon}

The photon sphere and critical impact parameter are most cleanly derived
in \emph{curvature} (Schwarzschild-like) coordinates, where the angular
part of the metric is manifestly $R^2\,\dd\Omega^2$ with $R$ the areal
radius.  We therefore begin by writing the TGP exponential metric in
this form.

\subsection{Exponential metric in curvature coordinates}

The isotropic-coordinate line element~\eqref{eq:exp-metric} is
equivalent, after a standard coordinate transformation, to the
\emph{Yilmaz--Rosen} form~\cite{Perlick2022,Rezzolla2014}:
\begin{equation}\label{eq:metric-curvature}
  \dd s^2 = -e^{-\rs/R}\,c_0^2\,\dd t^2
    + e^{+\rs/R}\,\dd R^2
    + R^2\,\dd\Omega^2,
\end{equation}
where $R$ is the areal radius and $\rs = 2GM/c_0^2$ as before.  We
define
\begin{equation}\label{eq:AB-def}
  A(R) \equiv e^{-\rs/R}, \qquad B(R) \equiv e^{+\rs/R}.
\end{equation}
Note that $A(R) > 0$ and $B(R) > 0$ for all $R > 0$: the metric is
regular everywhere outside the origin.

For comparison, the Schwarzschild metric has $A_{\mathrm{GR}} = 1 -
\rs/R$ and $B_{\mathrm{GR}} = (1 - \rs/R)^{-1}$, with $A_{\mathrm{GR}}
= 0$ at the horizon $R = \rs$.

\subsection{Null geodesics}

We restrict to the equatorial plane $\theta = \pi/2$ (spherical
symmetry).  The Killing vectors $\partial_t$ and $\partial_\varphi$
yield conserved energy~$E$ and angular momentum~$L$:
\begin{equation}\label{eq:conserved}
  E = A(R)\,c_0^2\,\dot{t},
  \qquad
  L = R^2\,\dot{\varphi},
\end{equation}
where the dot denotes differentiation with respect to an affine
parameter.  The null condition
$g_{\mu\nu}\dot{x}^\mu\dot{x}^\nu = 0$ gives the radial equation
\begin{equation}\label{eq:radial-null}
  \dot{R}^2 = \frac{E^2}{A(R)\,B(R)}
    - \frac{L^2}{R^2\,B(R)}
  = \frac{1}{B}\!\left(\frac{E^2}{A}
    - \frac{L^2}{R^2}\right).
\end{equation}
Since $A\,B = e^{-\rs/R}\cdot e^{+\rs/R} = 1$ for the exponential
metric, this simplifies to
\begin{equation}\label{eq:radial-simple}
  \dot{R}^2 = E^2\,e^{\rs/R}
    - \frac{L^2}{R^2}\,e^{\rs/R}
  = e^{\rs/R}\!\left(E^2 - \frac{L^2}{R^2}\right).
\end{equation}

Defining the impact parameter $b \equiv L/E$, the turning point
$\dot{R} = 0$ occurs where
\begin{equation}\label{eq:turning}
  b^2 = R^2.
\end{equation}
This would seem to imply that the closest approach is simply $R = b$
with no deflection.  However, this is an artifact of the product $AB = 1$
simplifying the radial equation.  The correct effective potential
analysis retains the full geodesic equation.

Let us instead use the standard photon-sphere criterion for a general
static spherically symmetric metric~\cite{Chandrasekhar1983,Perlick2022}.

\subsection{Photon sphere condition}\label{ssec:photon-sphere}

For a metric of the form
\begin{equation}
  \dd s^2 = -A(R)\,c_0^2\,\dd t^2 + B(R)\,\dd R^2 + R^2\,\dd\Omega^2,
\end{equation}
the effective potential for null geodesics, written in terms of the
impact parameter~$b$, takes the form
\begin{equation}\label{eq:Veff}
  V_{\mathrm{eff}}(R)
  = \frac{A(R)}{R^2}.
\end{equation}
The photon sphere is located where $V_{\mathrm{eff}}$ has an extremum:
\begin{equation}\label{eq:photon-sphere-areal}
  \frac{\dd}{\dd R}\!\left[\frac{A(R)}{R^2}\right] = 0
  \qquad\Longleftrightarrow\qquad
  \frac{A'(R)}{A(R)} = \frac{2}{R}.
\end{equation}

\paragraph{GR (Schwarzschild).}
With $A = 1 - \rs/R$:
\begin{equation}
  \frac{A'}{A} = \frac{\rs/R^2}{1 - \rs/R} = \frac{2}{R}
  \qquad\Longrightarrow\qquad
  R_{\mathrm{ph}}^{\mathrm{GR}} = \frac{3}{2}\,\rs.
\end{equation}

\paragraph{TGP (exponential).}
With $A = e^{-\rs/R}$:
\begin{equation}
  A'(R) = \frac{\rs}{R^2}\,e^{-\rs/R},
  \qquad
  \frac{A'}{A} = \frac{\rs}{R^2}.
\end{equation}
Setting this equal to $2/R$:
\begin{equation}\label{eq:rph-TGP}
  \frac{\rs}{R^2} = \frac{2}{R}
  \qquad\Longrightarrow\qquad
  \boxed{R_{\mathrm{ph}}^{\mathrm{TGP}} = \frac{\rs}{2}
  = \frac{GM}{c_0^2}.}
\end{equation}
The photon sphere in TGP lies at one-third the GR value: $R_{\mathrm{ph}}^{\mathrm{TGP}} = R_{\mathrm{ph}}^{\mathrm{GR}}/3$.

\subsection{Critical impact parameter}\label{ssec:bcrit}

The critical impact parameter is the value of~$b$ for a photon on the
unstable circular orbit at $R_{\mathrm{ph}}$:
\begin{equation}\label{eq:bcrit-formula}
  \bcrit = \frac{R_{\mathrm{ph}}}{\sqrt{A(R_{\mathrm{ph}})}}.
\end{equation}

\paragraph{GR (Schwarzschild).}
\begin{equation}\label{eq:bcrit-GR-full}
  \bcrit^{\mathrm{GR}}
  = \frac{3\rs/2}{\sqrt{1 - 2/3}}
  = \frac{3\rs/2}{\sqrt{1/3}}
  = \frac{3\sqrt{3}}{2}\,\rs
  = 3\sqrt{3}\,\frac{GM}{c_0^2}
  \approx 2.598\,\rs.
\end{equation}

\paragraph{TGP (exponential).}
\begin{equation}\label{eq:bcrit-TGP}
  \bcrit^{\mathrm{TGP}}
  = \frac{\rs/2}{\sqrt{e^{-\rs/(\rs/2)}}}
  = \frac{\rs/2}{\sqrt{e^{-2}}}
  = \frac{\rs}{2}\,e
  = \frac{GM\,e}{c_0^2}
  \approx 1.359\,\rs.
\end{equation}

The ratio of shadow sizes is therefore
\begin{equation}\label{eq:ratio}
  \boxed{
    \frac{\bcrit^{\mathrm{TGP}}}{\bcrit^{\mathrm{GR}}}
    = \frac{(\rs/2)\,e}{(3\sqrt{3}/2)\,\rs}
    = \frac{e}{3\sqrt{3}}
    \approx \frac{2.718}{5.196}
    \approx 0.523.
  }
\end{equation}
The TGP shadow is approximately $\mathbf{48\%}$ \textbf{smaller} than
the Schwarzschild prediction.

\begin{remark}[Relation to the $37\%$ estimate]\label{rem:37pct}
A preliminary estimate based on the isotropic-coordinate form of the
metric --- identifying the photon sphere with the isotropic coordinate
value $r_{\mathrm{ph}} = \rs$ and computing $\bcrit = \rs\sqrt{e}$ ---
yielded a ratio $\sqrt{e}/(3\sqrt{3}/2) \approx 0.635$, corresponding
to a $37\%$ reduction.  The discrepancy arises because the isotropic
coordinate~$r$ does not directly measure circumferential radius; the
physically meaningful calculation in curvature coordinates gives
$R_{\mathrm{ph}} = \rs/2$ and $\bcrit = (\rs/2)\,e$, producing the
larger $48\%$ reduction.  The qualitative conclusion --- that TGP
predicts a \emph{substantially} smaller shadow --- is the same in both
analyses.
\end{remark}

\subsection{Consistency check: does the exponential metric solve the TGP field equation?}
\label{ssec:consistency}

The preceding derivation assumed that the exponential metric holds at
all radii, including the photon sphere at $R = \rs/2$.  We now check
whether this assumption is self-consistent.

\paragraph{Residual of the field equation.}
The full TGP field equation in vacuum ($\rho = 0$) is
\begin{equation}\label{eq:field-eq-check}
  \psi'' + \frac{2}{r}\,\psi' + 2\,\frac{(\psi')^2}{\psi}
  + \beta\,\psi^2 - \gamma\,\psi^3 = 0,
\end{equation}
where $\beta = \gamma$ (vacuum condition).  For the exponential ansatz
$\psi_{\mathrm{exp}} = e^{-\rs/r}$, the kinetic part evaluates to
\begin{equation}
  \psi'' + \frac{2}{r}\psi' + \frac{2(\psi')^2}{\psi}
  = \frac{3\rs^2}{r^4}\,e^{-\rs/r},
\end{equation}
while the potential part is $\beta e^{-2\rs/r}(1 - e^{-\rs/r})$.  The
total residual
\begin{equation}\label{eq:residual}
  \mathcal{R}(r) = \frac{3\rs^2}{r^4}\,e^{-\rs/r}
  + \beta\,e^{-2\rs/r} - \gamma\,e^{-3\rs/r}
\end{equation}
is \textbf{strictly positive} for all $r > 0$ when $\beta = \gamma > 0$:
both the kinetic term ($> 0$) and the potential term
$\beta e^{-2\rs/r}(1-e^{-\rs/r}) > 0$ (since $e^{-\rs/r} < 1$) are
individually positive.

\begin{proposition}[Exponential is not the strong-field solution]
\label{prop:not-exact}
The exponential metric $\psi = e^{-\rs/r}$ does not satisfy the full
nonlinear TGP field equation~\eqref{eq:field-eq-check} at any radius.
It is an approximate solution valid only in the asymptotic regime
$r \gg \rs$ where $\mathcal{R}(r) \to 0$.
\end{proposition}

\paragraph{Absence of photon sphere in isotropic coordinates.}
The photon sphere condition in isotropic coordinates is
$\frac{d}{dr}[r\cdot\psi(r)] = 0$ (Section~\ref{ssec:photon-sphere}).
For the exponential metric:
\begin{equation}
  \frac{d}{dr}\!\left[r\,e^{-\rs/r}\right]
  = e^{-\rs/r}\!\left(1 + \frac{\rs}{r}\right) > 0
  \quad\text{for all } r > 0.
\end{equation}
The function $r\cdot\psi$ is \textbf{monotonically increasing}: there is
no extremum and hence \emph{no photon sphere} in coordinate-invariant
terms.  The result $R_{\mathrm{ph}} = \rs/2$ obtained in
Section~\ref{ssec:photon-sphere} arises from applying the curvature-coordinate
formula $A'(R)/A(R) = 2/R$ to the Yilmaz--Rosen form, but this form is
not the actual strong-field TGP solution (Proposition~\ref{prop:not-exact}).

\begin{remark}[Status of the $48\%$ result]\label{rem:48pct-status}
The $48\%$ shadow reduction derived in Section~\ref{ssec:bcrit} is a
\textbf{conditional result}: it holds \emph{if and only if} the
exponential metric is valid at $R = \rs/2$.
Proposition~\ref{prop:not-exact} shows this assumption is
inconsistent.  The actual TGP shadow size requires the full
nonlinear profile $\psi(r)$, derived in Section~\ref{sec:strong-field}.
\end{remark}

\subsection{Summary table: GR vs TGP}

\begin{table}[h]
\centering
\caption{Key black hole observables in GR (Schwarzschild)
  and TGP (exponential extrapolation, conditional ---
  see Proposition~\ref{prop:not-exact}).
  Areal radius $R$ throughout.}
\label{tab:comparison}
\begin{tabular}{lcc}
\toprule
Observable & Schwarzschild (GR) & TGP (exponential) \\
\midrule
Horizon radius $R_H$
  & $\rs = 2GM/c_0^2$
  & None (no horizon) \\[4pt]
Photon sphere $R_{\mathrm{ph}}$
  & $\frac{3}{2}\rs$
  & $\frac{1}{2}\rs$ \\[4pt]
Critical impact parameter $\bcrit$
  & $\frac{3\sqrt{3}}{2}\,\rs \approx 2.598\,\rs$
  & $\frac{e}{2}\,\rs \approx 1.359\,\rs$ \\[4pt]
Shadow ratio $\bcrit^{\mathrm{TGP}}/\bcrit^{\mathrm{GR}}$
  & $1$
  & $e/(3\sqrt{3}) \approx 0.523$ \\[4pt]
Curvature singularity
  & Yes ($R = 0$)
  & No ($K \to 0$) \\[4pt]
Event horizon
  & Yes
  & No (causal freezing) \\[4pt]
PPN parameters $\gamma, \beta$
  & $1, 1$
  & $1, 1$ \\
\bottomrule
\end{tabular}
\end{table}

% ============================================================
\section{No-singularity theorem}\label{sec:no-sing}

\subsection{Three mechanisms of singularity avoidance}

In TGP, the field-dependent constants provide three independent
mechanisms that prevent the formation of a curvature singularity:

\begin{enumerate}[(1)]
\item \textbf{Self-regulated gravity:} $G(\Phi) = G_0\,\PhiZero/\Phi
  \to 0$ as $\Phi \to \infty$.  The gravitational coupling itself
  vanishes in the strong-field interior, preventing further collapse.

\item \textbf{Causal freezing:} $c(\Phi) = c_0\sqrt{\PhiZero/\Phi}
  \to 0$ as $\Phi \to \infty$.  Light cones close continuously; no
  signal can propagate into the deep interior.  The region is
  \emph{causally frozen} --- operationally indistinguishable from a
  horizon, but without the global causal structure of a true event
  horizon.

\item \textbf{Quantum suppression:} $\hbar(\Phi) = \hbar_0\sqrt{%
  \PhiZero/\Phi} \to 0$ as $\Phi \to \infty$.  Quantum fluctuations,
  which in standard GR are expected to become important near a
  singularity, are \emph{suppressed} rather than enhanced.
\end{enumerate}

The combination of these three effects produces a ``frozen star''
interior: matter does not collapse to a point but is trapped in a
region of extreme field density where all dynamical processes cease.

\subsection{Kretschner scalar}

For the exponential metric~\eqref{eq:metric-curvature}, the Kretschner
scalar (the square of the Riemann tensor) can be computed directly.
Introducing the dimensionless variable $\chi \equiv \Phi/\PhiZero$ and
comparing with the Schwarzschild result $K_{\mathrm{GR}}$ at the same
coordinate radius:
\begin{equation}\label{eq:kretschner}
  K_{\mathrm{TGP}}(\chi) = \frac{K_{\mathrm{GR}}}{\chi^4}.
\end{equation}
As $\chi \to \infty$ (the strong-field interior), $K_{\mathrm{TGP}}
\to 0$: curvature \emph{decreases} rather than diverging.  This is the
opposite of the GR behavior and is a direct consequence of the
exponential metric structure.

\begin{proposition}[No-singularity theorem]\label{prop:no-sing}
In the TGP exponential metric, all polynomial curvature invariants
remain bounded for $r > 0$.  In particular,
$K_{\mathrm{TGP}}(r) \to 0$ as $r \to 0$.
\end{proposition}

\begin{proof}
The Kretschner scalar for the metric~\eqref{eq:metric-curvature} with
$\phi = -\rs/(2R)$ is a sum of terms proportional to
$e^{-n\rs/R}\cdot(\rs/R)^k$ for various positive $n$ and~$k$.
As $R \to 0$, each such term vanishes because $e^{-n\rs/R}$ decays
faster than any power of $R$.  Hence $K \to 0$.
\end{proof}

\subsection{Duality $\chi \leftrightarrow 1/\chi$}

The TGP framework possesses an \emph{operational duality} under the
inversion $\chi \mapsto 1/\chi$~\cite{Serafin2026-full}.  Under this
map:
\begin{itemize}[nosep]
  \item The cosmological vacuum ($\chi \to 0$, absence of spacetime)
    maps to the black hole interior ($\chi \to \infty$, frozen
    spacetime).
  \item In both limits, the information capacity $\mathcal{I} = 0$:
    no signal can propagate.
  \item The Planck length is invariant: $\lP \mapsto \lP$.
\end{itemize}
This duality is \emph{operational} (informational), not a full symmetry
of the field equations.  Both boundaries of the existence sector
$\mathcal{S}_1$ --- the ``nothingness from deficit'' ($\chi = 0$) and
the ``nothingness from excess'' ($\chi \to \infty$) --- are
functionally equivalent from the standpoint of an external observer.

The physical picture is that of a \textbf{frozen star}: the interior is
not empty (no singularity) but is a region of extremely dense substrate
field in which all physical processes have ceased.  This is the TGP
analog of the ``gravastar'' or ``frozen star'' proposals in the
literature, but it arises naturally from the field-dependent constants
rather than being imposed by hand.

% ============================================================
\section{Black hole thermodynamics in TGP}\label{sec:thermo}

A crucial consistency check for any modified gravity theory is whether
black hole thermodynamics survives.  In TGP, the Bekenstein--Hawking
entropy formula holds \emph{exactly}, with a remarkable cancellation
rooted in the Planck-length invariance.

\subsection{Bekenstein--Hawking entropy}

The entropy of a compact object with horizon-area~$A$ and local field
value $\Phi_H$ on its surface is
\begin{equation}\label{eq:S-BH-local}
  S_{\mathrm{BH}}
  = \frac{k_B\,c(\Phi_H)^3}{4\,G(\Phi_H)\,\hbar(\Phi_H)}\;A.
\end{equation}
Substituting the TGP scalings~\eqref{eq:scalings}:
\begin{equation}\label{eq:S-cancellation}
  \frac{c^3}{G\,\hbar}
  = \frac{c_0^3(\PhiZero/\Phi_H)^{3/2}}
         {G_0(\PhiZero/\Phi_H)\cdot\hbar_0(\PhiZero/\Phi_H)^{1/2}}
  = \frac{c_0^3}{G_0\,\hbar_0}\,
    \frac{(\PhiZero/\Phi_H)^{3/2}}{(\PhiZero/\Phi_H)^{3/2}}
  = \frac{c_0^3}{G_0\,\hbar_0}.
\end{equation}
The exponents cancel exactly: $\frac{3}{2} = 1 + \frac{1}{2}$.  This is
a direct consequence of $\lP = \sqrt{\hbar G/c^3} = \mathrm{const}$.
Therefore:
\begin{equation}\label{eq:S-BH-final}
  \boxed{S_{\mathrm{BH}}
  = \frac{k_B\,A}{4\,\lP^2},}
\end{equation}
which is the standard Bekenstein--Hawking formula expressed in
vacuum constants --- \textbf{completely independent of the local field
$\Phi_H$}~\cite{Bekenstein1973,Hawking1975}.

\subsection{Hawking temperature}

The Hawking temperature in TGP is modified by the local value of
$\hbar$ at the horizon~\cite{Serafin2026-full}:
\begin{equation}\label{eq:T-Hawking}
  T_H^{\mathrm{TGP}}
  = T_H^{\mathrm{GR}} \cdot \sqrt{\frac{\PhiZero}{\Phi_H}}
  \ll T_H^{\mathrm{GR}},
\end{equation}
since $\Phi_H \gg \PhiZero$ near the horizon.  The suppression of
Hawking radiation is consistent with the frozen-star picture: quantum
effects ($\hbar \to 0$) are suppressed where the field is strongest.

\subsection{First law of thermodynamics}

The first law $\dd E = T_H\,\dd S_{\mathrm{BH}}$ is preserved.  Since
$S_{\mathrm{BH}}$ is $\Phi$-independent (Eq.~\ref{eq:S-BH-final}) and
depends only on the ADM mass through $A = 16\pi G_0^2 M^2/c_0^4$
(at fixed $\Phi_H$), differentiation with respect to~$M$ reproduces the
standard thermodynamic identity~\cite{Bardeen1973}.

\begin{remark}[Consistency with Wald entropy]\label{rem:wald-entropy}
The Wald entropy formula~\cite{Wald1993} for diffeomorphism-invariant
theories gives $S = -2\pi\oint (\partial\mathcal{L}/\partial
R_{\mu\nu\rho\sigma})\,\hat{\epsilon}_{\mu\nu}\hat{\epsilon}_{\rho\sigma}
\,\dd\Sigma$.  In TGP, the effective gravitational Lagrangian inherits the
Einstein--Hilbert form with field-dependent $G$; the Wald calculation then
reproduces~\eqref{eq:S-BH-final} because the $\Phi$-dependence enters
only through $G(\Phi)$ and cancels in the entropy formula by the same
mechanism as~\eqref{eq:S-cancellation}.
\end{remark}

% ============================================================
\section{Observational predictions}\label{sec:obs}

\subsection{Shadow sizes for M87* and Sgr~A*}

Table~\ref{tab:shadows} compiles the predicted shadow angular diameters
for the two EHT targets.  The GR predictions use the Schwarzschild
metric (spin effects are subdominant for the shadow size at current
precision~\cite{EHT-M87-VI}).

\begin{table}[h]
\centering
\caption{Shadow angular diameters for M87* and Sgr~A* in GR
  and TGP (exponential extrapolation).  The TGP values are
  \textbf{conditional}: they assume the exponential metric holds at
  $R = \rs/2$, which Proposition~\ref{prop:not-exact} shows is
  inconsistent.  The true TGP prediction requires the full nonlinear
  profile (Section~\ref{sec:strong-field}).
  Observed values from~\cite{EHT-M87-VI,EHT-SgrA-2022}.}
\label{tab:shadows}
\begin{tabular}{lccc}
\toprule
Source & GR prediction & TGP prediction & Observed \\
\midrule
M87*
  & $\sim 42\;\mu\mathrm{as}$
  & $\sim 22\;\mu\mathrm{as}$
  & $42 \pm 3\;\mu\mathrm{as}$ \\[3pt]
Sgr~A*
  & $\sim 53\;\mu\mathrm{as}$
  & $\sim 28\;\mu\mathrm{as}$
  & $48.7 \pm 7\;\mu\mathrm{as}$ \\
\bottomrule
\end{tabular}
\end{table}

\subsection{Tension with EHT data}

The TGP prediction of a $\sim 48\%$ smaller shadow is in
\textbf{significant tension} with both EHT measurements.  For M87*, the
predicted $\sim 22\;\mu\mathrm{as}$ lies well outside the $1\sigma$
error bar of $42 \pm 3\;\mu\mathrm{as}$.  For Sgr~A*, the predicted
$\sim 28\;\mu\mathrm{as}$ is likewise incompatible with the observed
$48.7 \pm 7\;\mu\mathrm{as}$.

\subsection{Critical caveat: strong-field extrapolation}
\label{ssec:caveat}

\textbf{The $48\%$ figure rests on the assumption that the exponential
metric~\eqref{eq:metric-curvature} holds unchanged all the way down to
the photon sphere at $R = \rs/2$.}  This is a strong assumption.  In TGP, the exponential form is
derived as the unique conformal metric consistent with the field equation
in the \emph{weak-to-moderate} field regime.  The full radial profile
$\psi(r)$ in the strong-field regime --- where the nonlinear
self-interaction of $\Phi$ becomes important --- is an \textbf{open
problem}.

Several factors could moderate the discrepancy:
\begin{enumerate}[(a)]
\item \textbf{Nonlinear corrections.}\quad
  The TGP field equation contains cubic and quartic self-interaction
  terms ($\beta\chi^2 - \gamma\chi^3$) that are negligible at large~$r$
  but could modify $\psi(r)$ significantly for $r \lesssim \rs$.

\item \textbf{Interior structure.}\quad
  Unlike Schwarzschild, the TGP compact object has a finite interior
  (the frozen core).  The transition from the exterior exponential
  metric to the interior could shift the effective photon sphere
  outward.

\item \textbf{Accretion morphology.}\quad
  The EHT ``shadow'' is not a direct measurement of $\bcrit$ but of
  the bright emission ring, which depends on the accretion flow
  geometry.  In a modified metric, the ring-to-shadow mapping could
  differ.
\end{enumerate}

\textbf{We therefore present the $48\%$ reduction as an upper bound on
the deviation}, obtained from the simplest extrapolation of the TGP
metric.  The true deviation, once the full strong-field solution is
known, is expected to be smaller.  Even a $10\text{--}20\%$ deviation
would be a powerful discriminator between TGP and GR with next-generation
EHT data~\cite{Perlick2022}.

\subsection{Gravitational-wave predictions}

Independent of the shadow measurement, TGP predicts modifications to the
ringdown spectrum of merging compact objects:
\begin{itemize}[nosep]
  \item \textbf{Breathing mode} ($\ell = 0$):
    frequency shift $\delta f_0 \approx +0.6\%$ relative to
    GR~\cite{Serafin2026-full}.
  \item \textbf{Quadrupole mode} ($\ell = 2$):
    frequency shift $\delta f_2 \approx -0.2\%$.
\end{itemize}
These shifts are below the sensitivity of current LIGO/Virgo detectors
but are within reach of third-generation instruments (Einstein Telescope,
Cosmic Explorer) and of LISA for supermassive binary
mergers~\cite{Yunes2013,Yunes2009}.

Additionally, the scalar nature of the TGP field predicts a
\emph{breathing polarization} in the stochastic gravitational-wave
background, with a characteristic strain $h_{\mathrm{br}} \sim
10^{-17}$ at nanohertz frequencies --- potentially detectable by the
Square Kilometre Array (SKA) through pulsar timing
arrays~\cite{Yunes2013}.

% ============================================================
\section{Strong-field program}\label{sec:strong-field}

The consistency analysis of Section~\ref{ssec:consistency} demonstrates
that the exponential metric is insufficient in the strong-field regime.
We now outline what the correct strong-field solution must look like
and what it implies for the shadow.

\subsection{Boundary conditions from TGP ontology}

The TGP framework requires the substrate field to satisfy:
\begin{enumerate}[(i)]
\item $\psi(r \to \infty) = 1$ (vacuum value);
\item $\psi(r \to 0) \to \infty$ (frozen interior: $\Phi \to \infty$,
  $c \to 0$, $\hbar \to 0$, $G \to 0$).
\end{enumerate}
Note the contrast with the exponential ansatz, where
$\psi_{\mathrm{exp}}(r \to 0) = e^{-\rs/r} \to 0$ --- the opposite
behavior.  The full nonlinear solution must therefore deviate
qualitatively from the exponential form at small~$r$.

\subsection{Existence of a photon sphere: topological argument}

\begin{proposition}[Photon sphere existence]\label{prop:ps-existence}
Any smooth profile $\psi(r)$ with $\psi(0) = +\infty$ and
$\psi(\infty) = 1$ must have at least one photon sphere.
\end{proposition}

\begin{proof}
The photon-sphere condition (Section~\ref{ssec:photon-sphere}) is
$\frac{d}{dr}[r\,\psi(r)] = 0$, i.e., $\psi(r) + r\,\psi'(r) = 0$.
Consider the function $h(r) \equiv r\cdot\psi(r)$.  At $r = 0$:
$h(0) = 0$ (since $r = 0$, even though $\psi \to \infty$, the product
$r\cdot\psi \to 0$ provided $\psi$ diverges slower than $1/r$, which
is guaranteed by the finite-mass condition $\int \rho\,dV < \infty$).
At $r \to \infty$: $h(r) \to r \cdot 1 = \infty$ (monotonically).

For $\psi$ divergent at the origin, $h(r)$ starts at $h(0) = 0$,
rises to some value, and must eventually grow linearly.  If the
profile has a region where $\psi(r)$ decreases faster than $1/r$
(which it must, to transition from $\psi \gg 1$ near the center to
$\psi = 1$ at infinity), then $h'(r) = \psi + r\psi' < 0$ in that
region, producing a local minimum of~$h$.  At this minimum,
$h'(r_{\mathrm{ph}}) = 0$: a photon sphere exists.
\end{proof}

\subsection{Shadow size: qualitative bounds}

The critical impact parameter at the photon sphere is
$\bcrit = r_{\mathrm{ph}}\cdot\psi(r_{\mathrm{ph}})$ (in isotropic
coordinates).  The shadow ratio relative to GR is
\begin{equation}
  \frac{\bcrit^{\mathrm{TGP}}}{\bcrit^{\mathrm{GR}}}
  = \frac{r_{\mathrm{ph}}\,\psi(r_{\mathrm{ph}})}{(3\sqrt{3}/2)\,\rs}.
\end{equation}
Without the full profile $\psi(r)$, we can establish qualitative bounds:
\begin{itemize}[nosep]
  \item \textbf{Lower bound ($48\%$ smaller):} the exponential
    extrapolation, shown to be inconsistent
    (Proposition~\ref{prop:not-exact}).
  \item \textbf{Upper bound ($0\%$ deviation):} if the nonlinear
    corrections restore GR-like behavior ($\psi \approx 1 - \rs/r$
    near the photon sphere).
  \item \textbf{Realistic expectation:} the frozen interior
    ($\psi \to \infty$) pushes the effective metric toward
    $g_{tt} \to 0$, mimicking a horizon.  The photon sphere of
    such a metric is expected to lie \emph{between} the exponential
    value ($\rs/2$) and the GR value ($3\rs/2$), giving a shadow
    deviation of $\sim 0\text{--}48\%$.
\end{itemize}

\subsection{Numerical program (ongoing work)}

The resolution requires solving the full nonlinear ODE
\begin{equation}\label{eq:full-ode-program}
  \psi'' + \frac{2}{r}\psi' + \frac{2(\psi')^2}{\psi}
  + \beta\psi^2 - \gamma\psi^3 = -\frac{q}{\PhiZero}\,\rho(r)
\end{equation}
with a source term $\rho(r)$ that models the compact object, and
boundary conditions $\psi(0) \to \infty$, $\psi(\infty) = 1$.

Key challenges identified in preliminary numerical work~\cite{TGP-repo}:
\begin{enumerate}[(a)]
  \item The cubic self-interaction ($\gamma\psi^3$) causes runaway
    growth in shooting methods; relaxation or spectral methods are
    needed.
  \item The boundary condition $\psi \to \infty$ at $r = 0$ requires
    a regularized variable (e.g., $u = 1/\psi$ with $u(0) = 0$).
  \item The oscillatory tail at large~$r$ (Bessel-like, with
    wavelength $\lambda \sim 1/\sqrt{3}$ in code units) demands
    careful outer-boundary treatment.
\end{enumerate}
A complete strong-field solver implementing these techniques is
available in the TGP repository~\cite{TGP-repo} and will be the
subject of a dedicated companion paper.

% ============================================================
\section{Discussion}\label{sec:discussion}

\subsection{Honest assessment}

The analysis in this paper has revealed a subtlety missed in earlier
treatments: the exponential metric $\psi = e^{-\rs/r}$ does not satisfy
the full TGP field equation (Proposition~\ref{prop:not-exact}) and,
when examined in coordinate-invariant terms, does not possess a photon
sphere at all.  The ``$48\%$'' figure derived in
Section~\ref{ssec:bcrit} is therefore a \textbf{conditional result} that
does not represent the true TGP prediction.

This does not falsify TGP.  The situation is analogous to using the
Newtonian potential $U = -GM/r$ all the way to $r = 0$: the unphysical
result (infinite density) signals the breakdown of the approximation, not
of the underlying theory.  The exponential metric is the correct weak-field
solution; its extrapolation to $r \lesssim \rs$ is outside its domain of
validity.

The positive result (Section~\ref{sec:strong-field}) is that TGP's own
ontology --- the frozen interior with $\psi \to \infty$ --- guarantees
the existence of a photon sphere by a topological argument
(Proposition~\ref{prop:ps-existence}).  The shadow size is bounded between
$52\%$ (exponential limit) and $100\%$ (GR) of the Schwarzschild value,
with the exact answer depending on the full nonlinear profile.

\subsection{Toward the full strong-field solution}

The detailed program for obtaining the full nonlinear profile $\psi(r)$
is presented in Section~\ref{sec:strong-field}.  The key open question
is the shape of the transition region where $\psi(r)$ departs from the
weak-field exponential and begins to rise toward $\psi \to \infty$ at
the center.  This transition determines the photon sphere location and
hence the shadow size.  Until the full solution is obtained, TGP's
prediction for the BH shadow must be considered \textbf{undetermined}
rather than ``$48\%$ smaller.''

\subsection{Comparison with other modified gravity predictions}

Several modified gravity theories predict shadow sizes that deviate from
GR~\cite{Perlick2022,Cunha2018,Vagnozzi2023}:
\begin{itemize}[nosep]
  \item Scalar-tensor theories (Brans--Dicke): typically
    $\mathcal{O}(1\%)$ deviations for $\omega_{\mathrm{BD}} > 40000$.
  \item Einstein--dilaton--Gauss--Bonnet: up to $\sim 10\%$ for
    maximal coupling.
  \item Regular black holes (Bardeen, Hayward): $\sim 5\text{--}15\%$
    deviations depending on the regularization scale.
  \item Parametrized post-Einsteinian frameworks~\cite{Rezzolla2014,%
    Johannsen2013}: continuous deviation parameters.
\end{itemize}
The TGP prediction (even if moderated to $\sim 20\%$) would be among the
\emph{largest} shadow deviations predicted by any viable theory, making
it exceptionally testable.

\subsection{Distinguishing TGP from GR observationally}

Beyond the shadow size, TGP offers several independent observational
handles:
\begin{enumerate}[(i)]
\item \textbf{Shadow shape:}\quad
  The exponential metric predicts a different radial profile of
  deflection, leading to subtle differences in the photon ring
  structure (Lyapunov exponents, photon ring autocorrelation).

\item \textbf{Ringdown spectrum:}\quad
  The breathing mode shift $\delta f_0 \approx +0.6\%$ provides a
  gravitational-wave test independent of electromagnetic observations.

\item \textbf{Gravitational lensing:}\quad
  Strong lensing near the compact object will show deviations in the
  Einstein ring radius and time delays.

\item \textbf{Absence of information paradox:}\quad
  The frozen-star interior resolves the information paradox: information
  is not lost but trapped in a region of extremely slow (effectively
  zero) dynamics.  Late-time deviations in Hawking radiation spectrum
  could provide a test, though this is beyond current observational
  reach.
\end{enumerate}

% ============================================================
\section{Conclusion}\label{sec:conclusion}

We have analyzed the black hole shadow problem in TGP at three levels
of rigor, arriving at the following conclusions:

\begin{enumerate}[(1)]
  \item The weak-field exponential metric, if naively extrapolated,
    gives a photon sphere at $R_{\mathrm{ph}} = \rs/2$ and a shadow
    $48\%$ smaller than Schwarzschild.  However, this metric does
    \emph{not} satisfy the full TGP field equation
    (Proposition~\ref{prop:not-exact}) and possesses no photon sphere
    in coordinate-invariant terms.  \textbf{The $48\%$ figure is not
    the true TGP prediction.}

  \item The TGP ontology (frozen interior, $\psi \to \infty$) guarantees
    the existence of a photon sphere by a topological argument
    (Proposition~\ref{prop:ps-existence}).  The shadow size lies between
    $52\%$ and $100\%$ of the Schwarzschild value, with the exact answer
    requiring the full nonlinear profile $\psi(r)$.

  \item Three results are robust regardless of the strong-field details:
    (a)~TGP produces \emph{no curvature singularity}: $K \to 0$;
    (b)~black hole thermodynamics is preserved:
    $S_{\mathrm{BH}} = k_B A/(4\lP^2)$;
    (c)~the theory makes falsifiable predictions: shadow size (once
    computed), ringdown frequency shifts, and breathing-mode
    gravitational waves.
\end{enumerate}

The central open problem is the numerical solution of the full nonlinear
field equation with boundary conditions $\psi(0) \to \infty$,
$\psi(\infty) = 1$.  This will yield the definitive TGP prediction for
the BH shadow and determine whether TGP is consistent with EHT data.
Preliminary numerical infrastructure for this program is available
at~\cite{TGP-repo}.

TGP occupies a distinctive position: it is \emph{indistinguishable} from
GR in every weak-field test, yet its strong-field predictions ---
once fully computed --- will provide sharp, falsifiable targets for the
Event Horizon Telescope and gravitational-wave observatories.

\bigskip
\noindent\textbf{Acknowledgments.}\quad
This work was conducted as part of the TGP project.  Numerical tools and
derivations are available at~\cite{TGP-repo}.

% ============================================================
\bibliographystyle{unsrt}
\bibliography{tgp_bh_shadow}

\end{document}
